\chapter{Results and Discussion}
\label{chap:results}

The available experimental record establishes the cleaned sample, the
preprocessing audit, and the pre-thesis model comparison. It does not contain
the final output tables for SHAP rankings, fairness gaps, proxy prediction, or
reweighing. Those values must be generated by the pipeline described in
Chapter~\ref{chap:methodology}; they are not inferred here. Reporting them
without the experiment output would turn a methodological claim into an
unsupported numerical claim.

\section{Data Preparation Audit}

The audit found 165 exact duplicate rows and seven logically invalid records.
The earlier z-score rule deleted 1,700 observations and removed all applicants
aged 47 or older. It also deleted 7.1\% of the 25-and-above group, compared
with 2.1\% of the under-25 group. Finally, two engineered ratios recovered age
exactly and produced an age-predictability ROC-AUC of 1.000 in five-fold
cross-validation. These findings justify the revised preparation procedure and
show why a fairness result from the earlier pipeline would not be reliable.

\section{Pre-thesis Baseline Comparison}

\begin{table}[H]
\centering
\caption{Pre-thesis baseline comparison on the original protocol.}
\label{tab:baseline-results}
\begin{tabular}{lrrrrr}
\toprule
\textbf{Model} & \textbf{Accuracy} & \textbf{Precision} & \textbf{Recall}
& \textbf{F1} & \textbf{ROC-AUC}\\
\midrule
Logistic regression & 0.8569 & 0.7445 & 0.4896 & 0.5907 & 0.8632\\
Random forest & 0.8993 & 0.7705 & 0.7444 & 0.7572 & 0.9207\\
XGBoost & 0.9153 & 0.8085 & 0.7843 & 0.7963 & 0.9388\\
LightGBM & 0.9161 & 0.8054 & 0.7943 & 0.7998 & 0.9431\\
\bottomrule
\end{tabular}
\end{table}

LightGBM ranked first in ROC-AUC and recall in this preliminary comparison.
The result supports its use as the object of the thesis, but it does not
establish that LightGBM is universally superior to XGBoost: the comparison
uses one split and does not include a significance test. Its selection is also
motivated by the availability of exact tree explanations and the fact that the
research questions are defined for this model.

\section{Primary Research Results}

The following table is the required reporting structure for the final run.
It should be populated from the saved experiment output, using the same
primary split and the 95\% paired bootstrap intervals.

\begin{table}[H]
\centering
\caption{Required final comparison of the four model conditions.}
\label{tab:final-results}
\begin{tabular}{lrrrr}
\toprule
\textbf{Condition} & \textbf{ROC-AUC} & \textbf{Accuracy} & \textbf{DPD}
& \textbf{EOD}\\
\midrule
M0, age included & -- & -- & -- & --\\
M1, age removed & -- & -- & -- & --\\
M0-RW, reweighed & -- & -- & -- & --\\
M1-RW, reweighed & -- & -- & -- & --\\
\bottomrule
\end{tabular}
\end{table}

\section{Explainability and Proxy Findings}

The final report must list gain, split, and mean absolute SHAP rankings for
M0, together with Spearman correlation, Kendall's $\tau$, and top-k overlap.
For M1, it must report the age-predictability ROC-AUC and the SHAP gap
decomposition. Credit-history length, income, employment length, and loan
intent are the pre-specified candidates because they meet the association
screen. A candidate is not called a causal proxy unless its ablation changes
the measured gap and the evidence remains consistent across repeated splits.

\section{Mitigation Trade-off}

The final results must compare M0 with M0-RW and M1 with M1-RW. For every
comparison, the report should state the change in absolute DPD and EOD and
the corresponding change in ROC-AUC, accuracy, precision, recall, and F1.
A smaller gap without its performance cost is incomplete evidence; conversely,
a performance change without an uncertainty interval cannot establish a stable
mitigation effect.

\section{Limitations}

The dataset is simulated and may not represent a real lending portfolio. The
age distribution is concentrated near the young end, so conclusions about
older applicants are weak. The outcome is observed only for granted loans,
which creates a reject-inference limitation. Finally, demographic parity and
equal opportunity are only two fairness criteria, and SHAP is observational.
The final numerical conclusions should therefore be stated as findings about
this dataset, feature set, model, threshold, and experimental protocol.
